首先对附件的数据进行预处理，优化数据格式、处理缺失值和异常值，并且在考量测序质量后对相应数据优化和定权，以抑制低质量数据对后续问题分析的影响。

\subsection{问题一的分析}


问题一要求分析男胎Y染色体浓度与孕妇孕周数、BMI等指标的相关性，建立相应的关系模型并检验显著性。首先，考虑对某些孕妇有多次采血多次测量的情况，直接计算的相关系数于p值偏乐观，本文引入用\textbf{聚类自助法}计算斯皮尔曼相关系数的策略。建模阶段，先使用\textbf{线性混合模型LMM}，再进一步采用\textbf{广义加性混合模型GAMM}，通过似然比检验比较模型优劣，并输出显著性检验结果、效应方向和部分效应曲线，最后结合残差诊断和灵敏度分析验证模型稳健性，从而完成建模与求解。

\subsection{问题二的分析}

问题二聚焦于男胎孕妇的BMI对Y染色体浓度最早达标时间的影响，要求给出合理BMI分组和最佳检测时点，并分析检测误差对结果的影响。首先，利用线性插值和广义加性混合模型GAMM，针对不同情况估计孕妇\textbf{“最早达标时间”$T_i$}。再以加权负对数似然最小建立动态规划模型，同时结合BIC自动选择出组数$K^*$。再基于每组内$T_i$构造达标时间分布，定义\textbf{综合风险函数}，以风险最小为目标在合理搜索窗内求最优NIPT时点。最后把测量误差与个体差异纳入修正，得\textbf{误差修正}后的最优时点。最后做参数敏感性检查以分析检测误差对结果的影响。


\subsection{问题三的分析}

问题三进一步综合孕妇身高、体重、年龄等\textbf{多重因素}，同时考虑检测误差和浓度达标比例。首先进一步优化GAMM模型，综合多因素，估计个体在不同孕周下的\textbf{达标概率曲线}。然后以此为基础，利用\textbf{动态规划方法}在BMI轴上进行分段，结合综合风险函数求解合理的组数。最后在每个分组内确定能使风险函数最小化的最佳检测时点，从而得到\textbf{兼顾多因素差异与风险平衡}的分组和检测策略。

\subsection{问题四的分析}

问题四转向女胎的异常判定，关键是如何利用13号、18号、21号X染色体Z值、GC含量、读段数及比例等多维信息进行综合判断，实际上是一个\textbf{二分类问题}。首先为避免同一孕妇的多次检测数据被同时划入训练集和测试机，本文先按孕妇分组把数据训练集和测试集，随后做\textbf{特征工程}与\textbf{数据清洗}，筛选并标准化数值特征，得到核心候选特征组。对异常样本的\textbf{不平衡问题}，本文尝试构造多种\textbf{重采样}版本：随机过/欠采样、SMOTE、Borderline-SMOTE、SMOTEENN。然后基于特征数据，尝试两种\textbf{数集成模型}Gradient Boosting、AdaBoost与六种数据集分别搭配。最后均衡各项指标最佳选择总体最佳模型，利用\textbf{k折交叉验证}模型表现。最终可得判定异常女胎的模型方法$^{[8]}$。
